\section{Scope, Editions, and Qualifications}

\subsection{What this work is, and is not}

This is a study aid for reading the rubrics printed in front of the 1962
\work{Missale Romanum} and applying them to a particular day. It is not an
Ordo, not a liturgical book, not a permission, and not a substitute for a
competent superior's direction. It does not state, and nothing in it should be
read as stating, whether any person may lawfully celebrate with the 1962
Missal today: that is a question of present universal and particular law and
faculties, it is mutable, and it lies outside this study's scope and date.

Every worked answer is a derivation from the promulgated rubrics on stated
calendar facts. Where a result would depend on a national, diocesan,
religious, titular, patronal, dedication or indult entry, the result is stated
as a conditional and the approved proper is named as the thing that must be
produced. No particular calendar is supplied or invented here.

\subsection{Editions, and which text controls}

\begin{tabletwo}{Source}{Role in this work}
\work{Rubricae Breviarii et Missalis Romani}, as printed at the front of the
\work{Missale Romanum}, editio typica (Vatican City: Typis Polyglottis
Vaticanis, 1962): \work{Rubricae generales} nn.~1--137 and \work{Rubricae
generales Missalis Romani} nn.~269--530. & \textbf{Controlling witness} for
every rule cited as \rg{n} or \rgm{n}. Read in the Church Music Association of
America facsimile of that edition; the passages quoted or relied on were read
as page images at 260 dpi, not merely in the facsimile's optical text
layer.\\
The same code as promulgated in \work{Acta Apostolicae Sedis} 52 (1960)
597--705, with \work{Rubricarum instructum} at 593--595 and the Sacred
Congregation of Rites' general decree at 596. & \textbf{Textual control and
promulgation record.} Used to establish the promulgating act, its date and
effective date, and to supply \rgb{237}--\rgb{238}, which the Missal does not
print. Read in the Holy See's own scanned volume.\\
\work{Variationes in Breviario et Missali Romano ad normam novi codicis
rubricarum}, \work{Acta Apostolicae Sedis} 52 (1960) 706--740. & Used only for
what the reform abolished --- the seasonal orations and the Suffrage --- and
for nothing else.\\
The Missal's own \work{Ordo Missae}, \work{Ritus servandus} and per-Mass
rubrics, in the same facsimile. & \textbf{Controlling} for the text and
conditional directions of the Ordinary and for the directions printed with a
particular formulary.\\
Sacred Congregation of the Council, decree of 3 December 1960, \work{Acta
Apostolicae Sedis} 52 (1960) 985--986. & Cited once, in Section~10, for the
contemporaneous official sense of \latin{festa de praecepto}, expressly for the
\latin{pro populo} obligation and not as a rubrical definition.\\
\work{The Roman Missal translated into the English language for the use of the
laity} (Philadelphia: Eugene Cummiskey, 1861). & An identified nineteenth-century
English witness, quoted once. Pre-1955: its rubrics, seasonal second and third
orations, and Holy Week are not those of 1962.\\
Douay-Rheims Bible, Challoner revision. & An identified historical English
translation of the Vulgate, quoted once for the Last Gospel. Psalm numbering
throughout this work follows the Vulgate series the Missal cites.\\
\end{tabletwo}

The Benziger Brothers \work{editio iuxta typicam} of 1962 is a distinct
edition and was used only as an uncorrected optical-character transcription for
locating passages. It controls nothing here.

\subsection{Method, and how claims are graded}

Latin is quoted from the controlling witness. Ligatured \ae{} and \oe{} are
rendered \emph{ae} and \emph{oe}; the Missal's accent marks, which guide
pronunciation rather than sense, are not reproduced; nothing else is altered,
expanded or corrected. Where this book renders a rule in English it is an
explanation of the Latin, not a translation of it, and the Latin stands beside
it wherever the wording carries the argument.

Three grades of claim are used and are visible in the prose:

\begin{itemize}
\item \textbf{Direct rule} --- what a numbered rubric says, quoted or closely
restated at its own locus.
\item \textbf{Necessary inference} --- a result that follows from two or more
direct rules and is marked as an inference where it is not obvious. The reach
of \rgm{422}'s octave when All Souls has been transferred is the one place a
worked case rests on an inference, and it says so.
\item \textbf{Contested} --- a question the rubrics do not settle. These are
collected in Section~10 with both readings at full strength, and no answer is
asserted where we do not think the text supplies one.
\end{itemize}

Civil dates were computed from the Gregorian Easter and cross-checked against
the code's own date rules --- \rg{18}, \rg{17}\,d, \rg{20}, \rg{80} --- before
being used. Each worked case states the anchor dates it depends on so the
arithmetic can be reproduced. No year was taken from an Ordo, and no result in
this book has been checked against a published Ordo for the year it names.

\subsection{Bounds and known gaps}

\begin{itemize}
\item \textbf{The Office is out of scope.} Parts one and three of the code are
treated; part two, \work{Rubricae generales Breviarii Romani}, is used only for
\rgb{237}--\rgb{238}, because \rgm{431}\,a depends on it. Nothing here should be
used to arrange the Divine Office.

\item \textbf{Ceremonial is out of scope.} The \work{Ritus servandus} is used
for what it says about the text and its conditions, not for how ministers
move. Solemn Mass, pontifical Mass, ordination, funeral, marriage, procession
and exposition ceremonies belong to their own books.

\item \textbf{Holy Week is treated only at its edges.} The restored Order of
Holy Week is presupposed as the 1962 Missal prints it; its own extensive
rubrics are cited where a general rule points at them and are not expounded.

\item \textbf{No formulary is reproduced.} This work explains where text comes
from. It is not a hand missal and prints no Mass.

\item \textbf{Not collated against an Ordo.} Every worked case is a derivation.
An Ordo is compiled by a competent authority for a diocese or institute and
resolves locally the questions Section~10 leaves open.

\item \textbf{Canon law is not entered.} The obligation of hearing Mass, the
application \latin{pro populo}, marriage law and funeral law are named only
where a rubric points at them.
\end{itemize}

\subsection{Rights}

The Latin quoted from the 1960 code and the 1962 Missal is official liturgical
and legislative text of the Holy See; it is quoted at its own loci for study
and is not project-created expression. The English of the 1861 Cummiskey hand
missal and of the Douay-Rheims Challoner revision are identified historical
translations, quoted as witnesses and not offered for liturgical use; neither
was composed, adapted or modernised by this project. The facsimile and scanned
volumes through which these texts were read are third-party digitisations whose
rights status is recorded in this leaf's source records; no page image or bulk
transcription is reproduced here.

\subsection{Review state}

This document has been internally checked against the sources named above:
every numbered rubric cited was read at its own locus, and the passages quoted
were read as page images. That is a source audit, not an independent
specialist, theological, canonical or liturgical review, and it is not an
ecclesiastical approval of any kind. No \latin{imprimatur} or \latin{nihil
obstat} is claimed or implied.
